\chapter{The Harris decomposition}
\label{ch:harris}

Conditioned on survival, the skeleton of a Galton--Watson tree is again a Galton--Watson
tree, for a reduced law, and the subtrees hanging off it are independent Galton--Watson trees
for a conjugate law. This chapter proves both halves and then their joint form.

\section{Conditioning at the root}

\begin{definition}[The conditioned laws]
\label{def:conditioned}
\uses{def:sampleMeasure}
\lean{BranchingProcess.survivalMeasure, BranchingProcess.bushMeasure,
BranchingProcess.skeletonDegree}
\leanok
Write $\bbP^{\mathrm{s}}$ for the field measure conditioned on survival and
$\bbP^{\mathrm{b}}$ for it conditioned on extinction. The \emph{skeleton degree} of a field
is the number of children of the root founding an infinite subtree.
\end{definition}

\begin{lemma}[The surviving-children count]
\label{lem:surviveWeight}
\uses{def:conditioned}
\lean{BranchingProcess.Offspring.surviveWeight, BranchingProcess.sampleMeasure_skeletonDegree,
BranchingProcess.Offspring.surviveWeight_zero}
\leanok
Unconditionally, the root has exactly $k$ surviving children with probability
\[
  \sum_{j}\theta_j\binom{j}{k}(1-q)^{k}q^{\,j-k},
\]
and the value at $k=0$ is $f(q)=q$.
\end{lemma}

\begin{proof}
\leanok
Splitting at the root turns the event that the root has $j$ children of which exactly those in
a set $S$ survive into a product event, of mass $\theta_j(1-q)^{\abs{S}}q^{\,j-\abs{S}}$.
Summing over the sets of size $k$ gives the binomial count.
\end{proof}

\begin{definition}[The reduced and the conjugate law]
\label{def:reduced}
\uses{lem:surviveWeight, lem:extinction-zero}
\lean{BranchingProcess.Offspring.reduced, BranchingProcess.Offspring.conjugate}
\leanok
The \emph{reduced law} is the Harris--Sevastyanov transform
$\tilde\theta_k=\bigl(\sum_j\theta_j\binom{j}{k}(1-q)^kq^{\,j-k}\bigr)/(1-q)$ for $k\geq1$ and
$\tilde\theta_0=0$; the \emph{conjugate law} is $\hat\theta_j=\theta_jq^{\,j}/q$. Both are
offspring laws with the same support bound.
\end{definition}

\begin{proof}
\leanok
The reduced weights sum to one by the binomial theorem at $(1-q)+q=1$ together with
$\tilde\theta_0=q$ of \cref{lem:surviveWeight}; the conjugate weights sum to one because
$f(q)=q$.
\end{proof}

\section{The two halves}

\begin{theorem}[The skeleton is a Galton--Watson tree]
\label{thm:harris-skeleton}
\uses{def:reduced, def:sampleMeasure}
\lean{BranchingProcess.skeletonTreeLaw_eq_treeLaw}
\leanok
If $J\leq N$ and $q<1$, the re-addressed skeleton of a sample conditioned on survival has the
law of the Galton--Watson tree of the reduced law.
\end{theorem}

\begin{proof}
\leanok
Re-address the skeleton so that the $m$-th surviving child of a skeleton vertex is reached by
the letter $m$. \Cref{lem:surviveWeight} gives the arity at the root; the branching property
of \cref{lem:branching} makes each surviving subtree an independent copy of the conditioned
process, so the arity at an arbitrary skeleton vertex is reduced-law distributed and
independent of the address being present. The containment events of the finite prefix-closed
sets of words form a $\pi$-system generating the $\sigma$-algebra on subtrees, and the two
laws agree on it.
\end{proof}

\begin{theorem}[The bushes are Galton--Watson trees]
\label{thm:harris-bushes}
\uses{def:reduced, def:sampleMeasure}
\lean{BranchingProcess.bushTreeLaw_eq_treeLaw}
\leanok
If $J\leq N$ and $q>0$, a sample conditioned on extinction has the law of the Galton--Watson
tree of the conjugate law.
\end{theorem}

\begin{proof}
\leanok
The same route with the dying subtrees constrained: conditioned on extinction the root has
$j$ children with probability $\theta_jq^{\,j-1}$, and its subtrees are independent copies of
the process conditioned to die.
\end{proof}

\section{The joint form}

\begin{definition}[The decoration of a skeleton vertex]
\label{def:decoration}
\uses{thm:harris-skeleton, thm:harris-bushes}
\lean{BranchingProcess.decorationEvent, BranchingProcess.decorationMass}
\leanok
The \emph{decoration} of a skeleton vertex is its offspring count $j$, its number $k$ of
surviving children, and its dying subtrees indexed by rank.
\end{definition}

\begin{theorem}[The joint Harris decomposition]
\label{thm:harris-joint}
\uses{def:decoration}
\lean{BranchingProcess.survivalMeasure_decorations,
BranchingProcess.survivalMeasure_decorations_treeLaw}
\leanok
Conditioned on survival, probe the skeleton on a finite prefix-closed set $F$ of addresses and
constrain the decoration of each probed vertex. The mass factorises into one weight per
vertex,
\[
  \prod_{u\in F}
  \frac{\theta_{j_u}\binom{j_u}{k_u}(1-q)^{k_u}q^{\,j_u-k_u}}{1-q}
  \prod_{m<j_u-k_u}\bbP_{\hat\theta}(\cT\in\mathcal{A}_{u,m}),
\]
each dying subtree entering through the Galton--Watson law of the conjugate law.
\end{theorem}

\begin{proof}
\leanok
The joint root step constrains the surviving and the dying subtrees of one vertex at once,
through a product event over the blocks of the splitting. Decomposing a prefix-closed finite
set of addresses at the root and iterating that step down the skeleton gives the product, and
\cref{thm:harris-bushes} turns each bush factor into the conjugate tree law.
\end{proof}
